% !TeX program = xelatex
% 第二章 2.3、2.4 完整純文字稿。
% 接在既有 2.1、2.2 之後；不含 chapter 或手動章節編號。
% 資料記錄與實驗重現方式另安排於 2.7。
% 引用鍵：macenski2022ros2、nvidiaisaacsim、tang2024gmpc。

\section{軟體架構與模擬環境}
\label{sec:software_simulation}

本研究以 ROS 2 建立感知、規劃、控制與命令執行之間的通訊架構，並利用物理模擬環境進行方法開發與系統驗證。本節介紹軟體環境、通訊方式與模擬場景，作為後續系統整合及控制方法之基礎。

\subsection{ROS 2 軟體環境與通訊架構}
\label{subsec:ros2_architecture}

機器人作業系統第二代（Robot Operating System 2, ROS 2）提供模組化的軟體架構，使不同功能能以節點形式運作，並透過訊息介面交換資料\cite{macenski2022ros2}。本研究採用 ROS 2 Jazzy 作為主要通訊環境，依功能配置環境感知、機器人狀態取得、參考生成、預測控制及命令處理等模組，使各項功能能夠分別開發，並於整合後形成閉迴路控制流程。

各模組透過發布與訂閱機制傳遞雷射掃描、機器人狀態、參考路徑及控制命令等資訊。感測與狀態資料中的時間戳及座標框架，用於對應量測取得時刻與機器人位姿，使不同來源的資訊能轉換至控制所需的共同表示。控制與分析程式主要以 Python 實作，並利用 NumPy 與 SciPy 進行矩陣運算及稀疏矩陣處理；控制參數則透過設定檔管理，以支援不同任務及比較條件的配置。

\subsection{Isaac Sim 物理模擬與場景建置}
\label{subsec:isaac_sim_environment}

本研究以 NVIDIA Isaac Sim 作為移動操作與視覺導引任務之主要物理模擬平台。Isaac Sim 提供機器人模型、剛體與關節模擬、接觸處理及感測器模擬等功能，並支援以 Python 建立場景與控制模擬流程\cite{nvidiaisaacsim}。研究於模擬環境中配置移動機械臂、操作物件與室內設施，使控制命令能透過物理更新作用於機器人，再由狀態回授形成閉迴路。

場景依研究任務配置牆面、門扇、桌面、抽屜與障礙物等元素。其中，門扇與抽屜分別以旋轉及移動關節描述其機構運動，並依各實驗條件設定碰撞幾何、摩擦及關節參數。相機影像與深度資料用於視覺目標定位，底盤位姿、關節狀態與接觸資訊則支援控制及結果分析；各項資訊的使用方式與假設於第四章說明。前期動態導航測試另使用 Gazebo Harmonic，相關結果依其實際模擬平台與配置分別呈現。

物理更新、感測更新與控制求解具有不同的執行週期，因此，實驗中分別記錄模擬時間與主機實際運算時間，並保存物理步長、控制週期及模擬執行方式，作為分析運動表現與計算成本之依據。

\section{導航、視覺與操作之系統整合}
\label{sec:system_integration}

本研究將導航、視覺感知與操作控制連結為閉迴路系統。導航模組提供底盤行進參考，視覺模組提供操作目標資訊，控制器則依機器人狀態、環境資訊及任務需求產生運動命令。本節說明各模組之間的資訊傳遞，以及全身控制與命令執行之整合設計。

\subsection{全域路徑規劃與導航參考}
\label{subsec:global_planning_navigation_reference}

全域路徑規劃模組依據環境地圖、機器人起始位置與任務目的地，建立底盤的行進路徑。控制端根據當前位姿取得機器人在路徑上的進度，再沿路徑擷取預測時域內的參考點，形成位置、航向與速度參考，供底盤幾何模型預測控制使用。底盤追蹤之建模以 Tang 等人\cite{tang2024gmpc}提出的 GMPC 為相關基礎，並依本研究之全向底盤運動能力配置輸入限制。

在整合架構中，導航目的地可設定為門前、桌側或指定任務區域，精確的操作接近位置則依後續視覺量測與機器人狀態更新。當環境中存在移動障礙物時，障礙位置與運動估測提供予控制端，於追蹤導航參考的同時考量避障限制。進入操作區域後，底盤參考還需配合手臂可達範圍與物件機構條件調整，使移動與操作需求能夠共同進入控制問題。

\subsection{視覺感知與操作目標資訊}
\label{subsec:visual_perception_operation_targets}

視覺模組利用相機影像與深度資訊取得操作目標的三維位置，並將量測結果轉換至控制所使用的座標框架。底盤相機主要支援區域觀測與接近導引，手腕相機則提供近距離的目標確認與定位。目標資訊保留來源相機、影像擷取時間及量測有效性，使參考生成模組能依觀測時刻的機器人位姿進行座標轉換，並檢查不同來源之間的一致性。

本研究以指定任務區域與物件類型作為先驗資訊，再由當次視覺量測更新操作目標。參考生成模組結合目標位置、物件幾何及機器人構型，建立底盤接近參考與末端執行器目標。對於開門等受機構限制的操作，還需利用物件機構資訊描述抓持後允許的運動。目標資料的有效性與可執行條件作為參考承接的依據，使視覺量測能夠連結至後續控制。

\subsection{單一全身控制器與命令執行架構}
\label{subsec:wholebody_controller_execution}

本研究以單一全身預測控制器作為系統整合方向，規劃於共同的預測架構中處理導航、接近與操作需求。控制器依任務階段調整追蹤目標、權重及限制條件：導航階段著重底盤路徑追蹤與手臂姿態維持，接近階段加入視覺生成的末端目標，抓持後則進一步納入物件機構約束。各階段持續承接機器人狀態與命令執行資訊，使任務切換能配合當下運動條件進行。

控制器共同求解底盤速度與手臂設定點變化率，再由命令執行端轉換為底盤速度命令及手臂關節位置設定點。底盤命令依執行配置經過速度、變化率與輪級運動限制，手臂則由位置控制器追蹤設定點。由於命令處理可能改變最佳化所得的輸入，整合設計亦將實際放行命令與尚未生效的在途命令納入後續預測，以維持預測模型與執行過程之間的對應關係。

整體資料流由環境觀測與狀態回授開始，經參考生成、預測求解及命令執行後，再取得更新狀態進行下一控制週期。任務階段依目標到達、抓持或物件狀態等條件推進，相關切換條件、增廣預測模型與控制求解方法將於第三章說明。
